\section{Введение}
\label{sec:Section0} \index{Section0}

\section*{Введение}
С каждым годом интернет-инфраструктура занимает всё более значимую позицию в функционировании государственных, коммерческих и персональных информационных систем.
Экспоненциальный рост объёмов передаваемого трафика, обусловленный распространением облачных сервисов, мультимедийных потоков, IoT-устройств и распределённых вычислительных кластеров, формирует беспрецедентную нагрузку на сетевые подсистемы.
В этих условиях традиционные методы поверхностного анализа заголовков пакетов (shallow packet inspection) оказываются недостаточными для обеспечения требуемого уровня наблюдаемости, безопасности и соответствия регуляторным нормам.
Необходимость перехода к технологиям глубокого анализа сетевого трафика (Deep Packet Inspection, DPI) определяется требованием извлечения семантических признаков на уровне приложений, реконструкции сессий и контекстной интерпретации потоков данных \cite{zhang2021nta}.

Технология DPI подразумевает полный анализ структуры пакета, включая заголовки канального, сетевого, транспортного и прикладного уровней модели OSI.
В отличие от методов, оперирующих исключительно метаданными потока, DPI-анализаторы осуществляют парсинг полезных данных, идентификацию протоколов (в том числе зашифрованных и обфусцированных), извлечение полей заголовков прикладного уровня и формирование структурированного представления передаваемой информации.
Типичный конвейер DPI состоит из этапов захвата кадров, дефрагментации, реконструкции TCP-потоков, диссекции протоколов и сохранения извлечённых полей в нормализованном формате.
Однако полученный формат хранения, оптимизированный для эффективного разбора и компактного размещения в памяти, часто оказывается непрактичным для последующего анализа, корреляции и интеграции с системами мониторинга.
Сырые структуры данных избыточны и требуют дополнительной трансформации перед использованием аналитическими модулями или системами класса SIEM \cite{pang2020dpi}.

Для решения указанной проблемы применяется событийно-ориентированная архитектура (Event-Driven Architecture, EDA), в которой промежуточным звеном между низкоуровневым парсером трафика и высокоуровневыми анализаторами выступает система событий.
Система событий представляет собой набор механизмов, оперирующих двумя базовыми сущностями: событиями и обработчиками.
Событие определяется как типизированный набор данных, фиксирующий факт возникновения определённого состояния или действия в сетевом стеке (например, ``установлено TLS-соединение'', ``обнаружен аномальный DNS-запрос'', ``превышен порог пропускной способности для потока'').
Обработчик выступает в роли подписчика, принимающего события заданного типа и выполняющего детерминированную логику: агрегацию метрик, фильтрацию, обогащение контекстом или передачу во внешние хранилища.
Ключевым требованием к такой системе является возможность гибкой настройки схемы событий, динамической регистрации обработчиков, обеспечения типобезопасности и предсказуемой производительности в условиях высокой нагрузки \cite{fowler2014microservices}.

Существующие промышленные и академические решения в области анализа трафика часто реализуют жёстко закодированные или скриптово-зависимые механизмы генерации событий, что ограничивает их расширяемость, увеличивает задержки при обработке и усложняет интеграцию в существующие инфраструктуры.
Отсутствие унифицированного, высокопроизводительного и типобезопасного программного компонента, специально спроектированного для встраивания в DPI-конвейеры, обуславливает актуальность настоящего исследования.
Целью работы является изучение, проектирование и разработка библиотеки, предоставляющей функционал расширяемой системы генерации событий, пригодной для интеграции в различные анализаторы трафика.
Для достижения цели решаются следующие задачи: анализ архитектурных паттернов современных DPI-систем и существующих библиотек событий;
формулировка функциональных и нефункциональных требований;
проектирование API и внутренней архитектуры с учётом требований к масштабируемости и низкой задержке;
реализация прототипа на языке C++; экспериментальная оценка производительности, устойчивости и расширяемости разработанного решения \cite{inlinac2021survey}.

Научная новизна работы заключается в разработке компонентной модели событийного движка, сочетающей статическую типобезопасность на этапе компиляции (через механизмы CRTP и constexpr-идентификацию) с динамической маршрутизацией и асинхронной диспетчеризацией, оптимизированной для NUMA-архитектур и lock-free очередей.
Практическая значимость определяется возможностью использования библиотеки в качестве унифицированного промежуточного слоя между парсерами трафика и аналитическими модулями, что сокращает время разработки новых детекторов угроз, упрощает поддержку правил и снижает операционные затраты на интеграцию в существующие SOC-инфраструктуры.

\newpage

